\documentclass[10pt,a4paper]{amsart}
\usepackage[margin=19mm]{geometry}
\usepackage{amsmath,amssymb,mathtools}
\usepackage[scr=rsfs,cal=euler]{mathalfa}
\usepackage{xfrac}
\usepackage[unicode,hidelinks,bookmarks=false]{hyperref}
\newcommand{\ie}{i.e.\ }
\newcommand{\cf}{cf.\ }
\newcommand{\resp}{respectively\ }
\newcommand{\Subsection}{\subsection}
\providecommand{\nobreakdash}{\nobreak\hbox{-}}
\setlength{\emergencystretch}{2em}
\multlinegap0pt
\usepackage{bm}
\usepackage{bbm}
\usepackage{dsfont}
\usepackage{xspace}
\usepackage{cancel}
\usepackage{mathtools}
\usepackage{cleveref}
\usepackage{amsthm}
\makeatletter
\@ifundefined{thm}{\newtheorem{thm}{Thm}}{}
\@ifundefined{lem}{\newtheorem{lem}[thm]{Lemma}}{}
\makeatother
\newcommand{\Th}{\operatorname{Th}}
\newcommand{\cdim}{\operatorname{cd}}
\newcommand{\lk}{\operatorname{lk}}
\title[High dimensional hyperbolic Coxeter groups that virtually fi]{High dimensional hyperbolic Coxeter groups that virtually fiber}
\author{Jean-Francois Lafont, Barry Minemyer, Gangotryi Sorcar, Matthew Stover, Joseph Wells}
\date{Source: arXiv:2502.12906v2 (CC BY 4.0)}
\begin{document}
\maketitle

\noindent\emph{Source and licence.} High dimensional hyperbolic Coxeter groups that virtually fiber. Version arXiv:2502.12906v2 (CC BY 4.0), distributed under the Creative Commons Attribution 4.0 International licence, as recorded in the arXiv licence metadata for this exact version and shown on the abstract page https://arxiv.org/abs/2502.12906v2. Full rights notice: licenses/arxiv-2502.12906-CC-BY-4.0.txt.\par\smallskip

\noindent\textit{Editorial scope.} Counted unit: the Lem whose source label is \texttt{lem:cdim-of-links} in the article (source0.tex lines 1407--1409, source label \texttt{lem:cdim-of-links}), delivered together with the article's own complete original proof of it (source0.tex lines 1411--1425, one \texttt{proof} environment of 15 lines). No mathematics is written, repaired or re-derived: statement and proof are the article's own text line for line. Required context retained uncounted: lem:links-cube-to-vertex (lines 754--756); lem:three-props (lines 1046--1053). Every definition, display and label of those lines is retained unchanged. Omitted from this package and not redistributed: 1--753, 757--1045, 1054--1406, 1410--1410, 1426--1683. Nothing inside a retained excerpt is omitted or altered: every retained body is delivered exactly as the article's lines are written, so no omission receipt of any kind accompanies this package. Citations: the retained excerpts carry no citation command at all, so no bibliography entry of the article is transcribed and no reference is redistributed. Reference adaptation: none was needed and none was made; every reference inside the delivered unit resolves inside it, so no unresolved reference or citation is produced. Printed slips kept literal: no printed word, symbol, hypothesis or number of the counted statement or of its proof was altered. Layout and class: the article's own class is \texttt{article} with packages including verbatim, amsfonts, amsmath, amssymb, amsthm, caption, cleveref, enumitem, geometry, graphicx, stix, tikz, tikz-cd, xcolor; the wrapper is the lane's pinned \texttt{amsart} editorial wrapper at 10pt on A4, which restates the theorem-like environments the retained text needs and adds the article's own macro definitions that the retained text needs (\texttt{\textbackslash Th}, \texttt{\textbackslash cdim}, \texttt{\textbackslash lk}), with extra wrapper packages (bm, bbm, dsfont, xspace, cancel, mathtools, cleveref). The delivered numbering is the unit's own. Assets: no retained span contains a figure environment, a graphics-inclusion command, a PDF-page inclusion, a table or a TikZ picture; no image file is copied into this package, the package declares no asset dependency and no image file is redistributed with it. Licence: version arXiv:2502.12906v2 of this article is distributed under the Creative Commons Attribution 4.0 International licence, as recorded in the arXiv licence metadata for this exact version and shown on the abstract page https://arxiv.org/abs/2502.12906v2. A full rights notice is delivered at licenses/arxiv-2502.12906-CC-BY-4.0.txt. Characters: every retained excerpt is pure ASCII, so the delivered engine, XeLaTeX with its default Latin Modern text font, reports no missing character.
\begin{lem}\label{lem:links-cube-to-vertex}
    Let $X$ be an arbitrary cube complex, $\square \subset X$ be a cube, and $v\in \square$ be a vertex of the cube. The cube $\square$ defines a simplex $\tau \subset \lk(v,X)$, and there is a combinatorial isomorphism $\lk(\square, X) \cong \lk(\tau, \lk(v,X))$.
\end{lem}

\begin{lem}\label{lem:three-props}
Given any simplex $\sigma \subset \Th_\alpha(X)$, the link $\lk(\sigma, \Th_\alpha(X))$ contains a pair of subcomplexes $\lk^\prime(\sigma, \Th_\alpha(X))$, $\lk^{\prime \prime}(\sigma, \Th_\alpha(X))$ with the following three properties:
\begin{enumerate}
\item $\lk(\sigma, \Th_\alpha(X))$ decomposes as a join $\lk^\prime(\sigma, \Th_\alpha(X))*\lk^{\prime \prime}(\sigma, \Th_\alpha(X))$;
\item if $\lk^\prime(\sigma, \Th_\alpha(X))$ is nonempty, then it is contractible;
\item the complex $\lk^{\prime \prime}(\sigma, \Th_\alpha(X))$ is homotopy equivalent to $\lk(\square_\sigma , X)$.
\end{enumerate}
\end{lem}

\begin{lem}\label{lem:cdim-of-links}
    If $\sigma$ is an arbitrary nonempty simplex in $T_n$, then $\cdim(\lk(\sigma, T_n)) \leq n-1$.
\end{lem}

\begin{proof}
Since $T_n$ is a thickening, we can apply \Cref{lem:three-props} to see that $\lk(\sigma, T_n)$ decomposes as a join $\lk^\prime(\sigma, T_n)*\lk^{\prime\prime}(\sigma, T_n)$. If $\lk^\prime(\sigma, T_n)\neq \emptyset$, then $\lk^\prime(\sigma, T_n)$ is a nonempty simplex, hence it is contractible. This forces $\lk(\sigma, T_n)$ to likewise be contractible, in which case we are done.

Alternatively, if $\lk^\prime(\sigma, T_n)= \emptyset$, then \Cref{lem:three-props} tells us that
    \[ \lk(\sigma, T_n) = \lk^{\prime\prime}(\sigma, T_n) \simeq \lk(\square_\sigma, X_n). \]
We let $v\in \square_\sigma$ be an arbitrary vertex of the cube $\square_\sigma$, and apply \Cref{lem:links-cube-to-vertex} to obtain a combinatorial isomorphism
    \[ \lk(\square_\sigma, X_n) \cong \lk(\tau, \lk(v,X)) \cong \lk(\tau, T_{n-1}) \]
where $\tau$ is the (possibly empty) simplex in $\lk(v,X)$ corresponding to $\square_\sigma$.

If the simplex $\tau$ is empty, this means $\square_\sigma$ coincides with the vertex $v$ and $\lk(\tau, T_{n-1}) = T_{n-1}$. In that case we conclude by induction with respect to proving the Main Theorem that
    \[ \cdim(\lk(\sigma, T_n)) = \cdim(T_{n-1}) = n-1. \]
On the other hand, if $\tau$ is nonempty then we again obtain
    \[ \cdim(\lk(\sigma, T_n)) = \cdim(\lk(\tau,T_{n-1})) \leq n-2 \]
by induction with respect to both this lemma and the Main Theorem. Either way we obtain the desired bound, which concludes the proof of the lemma.
\end{proof}

\end{document}
